\usemodule[memos]
\usemodule[basicexam]
\usemodule[setups]

\enablemode[no-setup-main]
\loadsetups[i-basicexam.xml]

\setupinteraction[state=start]
\setuptextrules[style=\tf\ss,before=,after=,rulethickness=.5pt]

\setuptyping[style=\ttxx\setuplocalinterlinespace[line=2ex],
             option={context},
             before={\vbox{\hrule height .5pt\par
                     \textrule{Code Example}%
                     \hrule height .5pt\vrule width0pt depth 4pt}\par},
              after={\vbox{\hrule height .5pt\strut}},
             ]

\starttext

\title{basicexam Module User Manual}

\blank[big]

\placecontent

\chapter{Module Overview}

\section{Introduction}

basicexam is an exam paper production module designed for ConTeXt LMTX, providing a complete exam typesetting solution.
The module supports various question types, including multiple choice, fill-in-the-blank, short answer, cloze, etc., and has features like answer management and score statistics.

\section{Main Features}

\startitemize[packed]
\item Multiple question types: multiple choice, fill-in-the-blank, short answer, cloze, etc.
\item Teacher/Student mode switching to control answer display
\item Unlimited nested question support
\item Automatic score statistics and answer collection
\item Flexible style customization options
\item Built-in various symbol markers
\item Composition grids and writing areas
\item Reading material and reference system
\stopitemize

\section{System Requirements}

This module requires {\bf ConTeXt LMTX (LuaMetaTeX)} environment and does not support MkIV version.
If running in an incompatible environment, the module will display an error and terminate compilation.

\chapter{Module Loading and Parameters}

\section{Loading the Module}

\starttyping
\usemodule[basicexam]
\stoptyping

\section{Module Parameters}

The following parameters can be set when loading the module:

\starttabulate[|l|l|p|]
\HL
\NC Parameter \NC Options \NC Description \NC\NR
\HL
\NC mode \NC \NC Running mode: \NC\NR
\NC      \NC student  \NC Student mode, hide answers \NC\NR
\NC      \NC teacher  \NC Teacher mode, show answers \NC\NR
\NC      \NC check    \NC Check mode, show answers and output debug info \NC\NR
\HL
\NC layout \NC \NC Layout settings: \NC\NR
\NC        \NC twoup   \NC Two-up layout (two pages per sheet) \NC\NR
\NC        \NC none    \NC Default single page layout \NC\NR
\HL
\NC footer \NC \NC Footer settings: \NC\NR
\NC        \NC default \NC Show subject and page number \NC\NR
\NC        \NC none    \NC No footer \NC\NR
\HL
\NC libpath \NC path \NC SQLite library file path (for data storage) \NC\NR
\HL
\NC showanswer \NC showanswer \NC Show answers \NC\NR
\NC showpoint  \NC showpoint  \NC Show scores \NC\NR
\NC showscript \NC showscript \NC Show answer area \NC\NR
\HL
\stoptabulate

\subsection{Examples}

\starttyping
% Teacher mode, show answers
\usemodule[basicexam][mode=teacher]

% Student mode, hide answers
\usemodule[basicexam][mode=student]

% Check mode, debug output
\usemodule[basicexam][mode={teacher,check}]

% Two-up layout
\usemodule[basicexam][layout=twoup]
\stoptyping

\section{Mode Control}

\subsection{Global Mode}

Set global mode through module parameters:

\starttyping
\usemodule[basicexam][mode=teacher]  % Teacher mode
\usemodule[basicexam][mode=student]  % Student mode
\stoptyping

\subsection{Local Control}

Dynamically switch in the document:

\starttyping
\showanswer   % Show answers
\hideanswer   % Hide answers
\showpoint    % Show scores
\hidepoint    % Hide scores
\stoptyping

\subsection{Individual Question Control}

Control display at question level:

\starttyping
\startquestion[showanswer=true,showpoint=true]
  This question will show answers and scores.
\stopquestion

\startquestion[showanswer=false,showpoint=false]
  This question will hide answers and scores.
\stopquestion
\stoptyping

\section{Style Customization}

\subsection{Global Settings}

Use \type{\setupquestion} for global settings:

\startbuffer
\setupquestion[
  answer={default answer},
  point=1,
  showanswer=true,
  showpoint=true,
  style=\ss,
  color=black,
]
\stopbuffer

\typebuffer

\subsection{Multiple Choice Style}

\startbuffer
\setupchoice[
  answerstyle=\ss,
  answercolor=red,
  numberconversion=A,
  option={A,packed,joinedup,nowhite},
]
\stopbuffer

\typebuffer

\subsection{Answer Style}

\startbuffer
\setupanswer[
  style=\ss\it,
  color=darkblue,
  answerlabel={{Answer}},
  pointlabel={, points},
  before=\blank[halfline]\startnarrower,
  after=\blank[halfline]\stopnarrower,
]
\stopbuffer

\typebuffer

\chapter{Paper Title}

\section{papertitle Command}

The paper title is an important part of an exam paper, usually containing exam name, subject, grade, exam time, etc.
The module provides \type{\definepapertitle} and \type{\makepapertitle} commands to create standardized paper titles.

\subsection{Basic Usage}

First use \type{\definepapertitle} to define title style, then use \type{\setuppapertitle} to set content, finally use \type{\makepapertitle} to output:

\startbuffer
\definepapertitle[myexam][list={secret,title,subject,information}]
\setuppapertitle[myexam][
  secret={CONFIDENTIAL $\star$ DO NOT OPEN},
  title={2024 National College Entrance Examination},
  subject={Mathematics},
  information={Total: 150 points, Time: 120 minutes}]
\makepapertitle[myexam]
\stopbuffer

\typebuffer

\getbuffer

\subsection{Title Elements}

The \type{list} parameter controls which elements are displayed in the title and their order. Default list: \type{secret,title,subject,information,signinformation,notice}.

\starttabulate[|l|p(8cm)|]
\HL
\NC Element \NC Description \NC\NR
\HL
\NC secret \NC Confidentiality mark (e.g., "CONFIDENTIAL $\star$ DO NOT OPEN") \NC\NR
\NC title \NC Main title \NC\NR
\NC subject \NC Subject name (automatically adds subjectsuffix) \NC\NR
\NC information \NC Exam information (e.g., total points, time) \NC\NR
\NC signinformation \NC Signature information area \NC\NR
\NC notice \NC Notice \NC\NR
\HL
\stoptabulate

Each element can be styled individually with style, color, and alignment, e.g.:
\type{secretstyle}, \type{secretcolor}, \type{secretalign}, etc.

\subsection{Other Parameters}

\starttabulate[|l|l|p|]
\HL
\NC Parameter \NC Default \NC Description \NC\NR
\HL
\NC alternative \NC default \NC Title template \NC\NR
\NC list \NC \NC Display element list \NC\NR
\NC subjectsuffix \NC Exam Paper \NC Subject suffix \NC\NR
\NC style \NC \NC Global font style \NC\NR
\NC color \NC \NC Global font color \NC\NR
\NC align \NC center \NC Global alignment \NC\NR
\HL
\stoptabulate

\section{Custom Title Style}

You can create new title styles by inheriting existing ones:

\startbuffer
\definepapertitle[mytitle][myexam]
\setuppapertitle[mytitle][title={My Exam Paper},style=\ssa\bf]
\makepapertitle[mytitle]
\stopbuffer

\typebuffer

\start\getbuffer\stop

\chapter{Multiple Choice}

\section{Basic Usage}

Multiple choice questions have two syntax formats: {\bf environment format} and {\bf shortcut format}.

\subsection{Environment Format}

Use \type{\startquestion}...\type{\stopquestion} and \type{\startchoice}...\type{\stopchoice} environments:

\startbuffer
\startquestion
  This is the question stem.
  \startchoice
    \startcitem Option A \stopcitem
    \startcitem Option B \stopcitem
    \startcitem Option C \stopcitem
    \startcitem Option D \stopcitem
  \stopchoice
\stopquestion
\stopbuffer

\typebuffer \getbuffer

\subsection{Shortcut Format}

Use \type{\question} and \type{\choice} commands:

\startbuffer
\question{This is the question stem.
  \choice{Option A,Option B,Option C,Option D}
}
\stopbuffer

\typebuffer \getbuffer

\section{Parameter Description}

\showsetup[setupchoice]

\starttabulate[|l|l|p|]
\HL
\NC Parameter \NC Default \NC Description \NC\NR
\HL
\NC inlinechoice \NC false \NC Whether it is an inline multiple choice \NC\NR
\NC inlineleft \NC \NC Content on the left of inline options \NC\NR
\NC inlineright \NC \NC Content on the right of inline options \NC\NR
\NC correctsymbol \NC \NC Correct answer marker symbol \NC\NR
\NC separator \NC \NC Option separator \NC\NR
\NC textstyle \NC \NC Option text style \NC\NR
\NC textcolor \NC \NC Option text color \NC\NR
\NC showanswer \NC \NC Whether to show answer \NC\NR
\NC showpoint \NC \NC Whether to show score \NC\NR
\NC answer \NC \NC Answer content \NC\NR
\NC answerstyle \NC \NC Answer style \NC\NR
\NC answercolor \NC \NC Answer color \NC\NR
\NC answerlabel \NC \NC Answer label \NC\NR
\NC point \NC \NC Question score \NC\NR
\NC pointstyle \NC \NC Score style \NC\NR
\NC pointcolor \NC \NC Score color \NC\NR
\NC pointlabel \NC \NC Score label \NC\NR
\HL
\stoptabulate

\section{Advanced Usage}

\subsection{Marking Correct Answers}

Use \type{[*]} to mark correct answers:

\startbuffer
\startquestion[showanswer=true]
  This is the question stem.
  \startchoice
    \startcitem    Option A \stopcitem
    \startcitem    Option B \stopcitem
    \startcitem[*] Option C \stopcitem
    \startcitem    Option D \stopcitem
  \stopchoice
\stopquestion

\question[showanswer=true]{This is the question stem.
  \choice{Option A,Option B,{[*]Option C},Option D}
}
\stopbuffer

\typebuffer \getbuffer

\subsection{Multiple Correct Answers}

Support for marking multiple correct answers:

\startbuffer
\startquestion[showanswer=true]
  Select all correct options.
  \startchoice
    \startcitem[*] Option A \stopcitem
    \startcitem    Option B \stopcitem
    \startcitem[*] Option C \stopcitem
    \startcitem    Option D \stopcitem
  \stopchoice
\stopquestion
\stopbuffer

\typebuffer \getbuffer

\subsection{Column Control}

Control the number of display columns through optional parameters:

\startbuffer
\startquestion
  Options displayed in 8 columns.
  \startchoice[n=8]
    \startcitem A \stopcitem
    \startcitem B \stopcitem
    \startcitem C \stopcitem
    \startcitem D \stopcitem
    \startcitem E \stopcitem
    \startcitem F \stopcitem
    \startcitem G \stopcitem
    \startcitem H \stopcitem
  \stopchoice
\stopquestion

\question{Options displayed in 8 columns.
  \choice[n=8]{A,B,C,D,E,F,G,H}
}
\stopbuffer

\typebuffer \getbuffer

\subsection{Inline Multiple Choice}

Inline multiple choice displays options within the question stem line, suitable for short questions:

\startbuffer
\startquestion
  This is an inline multiple choice.\relax
  \startinlinechoice
    \startcitem Option A \stopcitem
    \startcitem Option B \stopcitem
    \startcitem Option C \stopcitem
    \startcitem Option D \stopcitem
  \stopinlinechoice
\stopquestion

\question{This is an inline multiple choice.\relax
  \inlinechoice{Option A,Option B,Option C,Option D}
}
\stopbuffer

\typebuffer \getbuffer

\section{Custom Styles}

Use \type{\definechoice} to create custom multiple choice styles:

\startbuffer
\definechoice[mychoice][
  option={text,6,packed,joinedup,nowhite},
  inlinechoice=true,
  correctsymbol={\symbol{markcirclecheck}},
  inlineleft={【\nobreak},
  inlineright={\nobreak 】},
  answerstyle=\ss,
  separator={\removeunwantedspaces,\space},
]

\startquestion[showanswer=true]
  Using custom multiple choice style.
  \startmychoice
    \startcitem[*] Option A \stopcitem
    \startcitem    Option B \stopcitem
    \startcitem    Option C \stopcitem
    \startcitem    Option D \stopcitem
  \stopmychoice
\stopquestion
\stopbuffer

\typebuffer \getbuffer

\chapter{Questions and Problems}

\section{question Environment}

\type{question} environment is used to create independent questions:

\startbuffer
\startquestion[option={n,packed,joinedup,nowhite}]
This is the main content of the question.
\stopquestion

\question{This is the main content of the question.}
\stopbuffer

\typebuffer \getbuffer

\section{problem Environment}

\type{problem} environment is used to create sub-problem lists:

\startbuffer
\startquestion
This is the stem of a major question.
  \startproblem
  \startpitem This is the first sub-question.\stoppitem
  \startpitem This is the second sub-question.\stoppitem
  \stopproblem
\stopquestion

\question{This is the stem of a major question.
  \problem[textstyle=ss]
    {\pitem{This is the first sub-question.}
     \pitem{This is the second sub-question.}
  }
}
\stopbuffer

\typebuffer \getbuffer

\section{Question Nesting}

Support for unlimited levels of question nesting:

\startbuffer
\startquestion
This is the first level question.
  \startquestion
  This is the second level question.
    \startquestion
    This is the third level question.
    \stopquestion
  \stopquestion
\stopquestion
\stopbuffer

\typebuffer \getbuffer

\section{question Parameters}

\showsetup[setupquestion]

\chapter{Answers and Analysis}

\section{Basic Usage}

\type{answer} environment is used to add answers and analysis:

\startbuffer
\startquestion[point=4,showpoint=true,answer=A]
This is the question stem.
  \startchoice
    \startcitem[*] Option A \stopcitem
    \startcitem    Option B \stopcitem
    \startcitem    Option C \stopcitem
    \startcitem    Option D \stopcitem
  \stopchoice
  \startanswer
    This is a detailed explanation of the question.
    The reason for choosing A is...
  \stopanswer
\stopquestion
\stopbuffer

\typebuffer \getbuffer

\section{Parameter Description}

\starttabulate[|l|l|p|]
\HL
\NC Parameter \NC Default \NC Description \NC\NR
\HL
\NC showanswer \NC \NC Whether to show answer \NC\NR
\NC showpoint \NC \NC Whether to show score \NC\NR
\NC showscript \NC \NC Whether to show answer area \NC\NR
\NC scriptcontent \NC \NC Answer area prompt content \NC\NR
\NC style \NC \ss\it \NC Answer font style \NC\NR
\NC color \NC \NC Answer color \NC\NR
\NC answerlabel \NC Answer \NC Answer label \NC\NR
\NC pointlabel \NC , points \NC Score label \NC\NR
\HL
\stoptabulate

\section{Advanced Usage}

\subsection{Answer Area}

Use \type{showscript=true} to add answer writing area:

\startbuffer
\startquestion[point=4,showpoint=true]
This is a question that requires a written answer.
\startanswer[showscript=true,frame=off,
             scriptcontent={\centerline{\lightgray\ss Please write your answer here}\par
                            \null\vskip3\baselineskip\relax}]
\stopanswer
\stopquestion
\stopbuffer

\typebuffer \getbuffer

\subsection{Score Settings}

Set scores in questions:

\startbuffer
\question[point=5,showpoint=true]{This question is worth 5 points.\choice{A,B,C,D}}

\startquestion[point=10,showpoint=true]
This major question is worth 10 points total.
\startproblem
  \startpitem[point=3,showpoint=true] First sub-question is 3 points.\stoppitem
  \startpitem[point=4,showpoint=true] Second sub-question is 4 points.\stoppitem
  \startpitem[point=3,showpoint=true] Third sub-question is 3 points.\stoppitem
\stopproblem
\stopquestion
\stopbuffer

\typebuffer \getbuffer

\subsection{score Command}

\type{\score} command is used to mark the score for each point in the answer, helping graders understand score distribution.
This command is usually used with \type{\resetscore}, which resets the cumulative score.

Use \type{\score{points}} in answers to mark scoring points:

\startbuffer
\startanswer
  Answer point one \score{3}
  Answer point two \score{4}
  Answer point three \score{3}
  \resetscore
  Total \score[type=dotfills]{10}
\stopanswer
\stopbuffer

\typebuffer \getbuffer

score command parameters:

\starttabulate[|l|l|p|]
\HL
\NC Parameter \NC Default \NC Description \NC\NR
\HL
\NC type \NC \NC Display type: dotfills, normal \NC\NR
\NC style \NC \NC Number style \NC\NR
\NC color \NC \NC Number color \NC\NR
\NC before \NC \NC Content before score \NC\NR
\NC after \NC \NC Content after score \NC\NR
\HL
\stoptabulate

\subsection{Answer Output}

Answer output function is used to generate answer summary tables at the end of the paper or at specified locations.
Use \type{\typeanswer} command to output answers in different formats.

\startbuffer
% List format (default)
\typeanswer[start=1,total=10]

% Table format
\typeanswer[alternative=table,start=1,total=30,rows=6,columns=5]

% Grid format
\typeanswer[alternative=grid,start=1,total=10,rows=5,columns=2]

% Direct access to individual answer
\typeanswer[alternative=direct]{Q-5}
\stopbuffer

\typebuffer

typeanswer command parameters:

\starttabulate[|l|l|p|]
\HL
\NC Parameter \NC Default \NC Description \NC\NR
\HL
\NC alternative \NC list \NC Output format: list, table, grid, direct \NC\NR
\NC start \NC 1 \NC Starting question number \NC\NR
\NC stop \NC \NC Ending question number \NC\NR
\NC total \NC \NC Total number of questions (alternative to stop) \NC\NR
\NC chapter \NC \NC Chapter number \NC\NR
\NC rows \NC \NC Number of table rows \NC\NR
\NC columns \NC \NC Number of table columns \NC\NR
\NC separator \NC \NC Separator (text format) \NC\NR
\NC left \NC \NC Content to the left of answer \NC\NR
\NC right \NC \NC Content to the right of answer \NC\NR
\HL
\stoptabulate

\chapter{Fill-in-the-Blank}

\section{Basic Usage}

\type{\fillin} command is used to create fill-in-the-blank items in questions:

\startbuffer
\startquestion[showanswer=true,point=2,showpoint=true]
  Lorem ipsum dolor sit amet,
  \fillin{fill-in answer}
  consectetur adipiscing elit.
\stopquestion
\stopbuffer

\typebuffer \getbuffer

\section{Parameter Description}

\showsetup[setupfillin]

\starttabulate[|l|l|p|]
\HL
\NC Parameter \NC Default \NC Description \NC\NR
\HL
\NC style \NC \NC Fill-in content style \NC\NR
\NC color \NC \NC Fill-in content color \NC\NR
\NC before \NC \NC Content before fill-in \NC\NR
\NC after \NC \NC Content after fill-in \NC\NR
\NC left \NC \NC Left boundary \NC\NR
\NC right \NC \NC Right boundary \NC\NR
\NC empty \NC yes \NC Blank display mode: yes, number, none \NC\NR
\NC alternative \NC default \NC Fill-in style: default, frame, simple \NC\NR
\NC n \NC \NC Underline length (number of characters) \NC\NR
\NC width \NC \NC Fill-in width \NC\NR
\HL
\stoptabulate

\section{Advanced Usage}

\subsection{Fill-in Styles}

Control fill-in display through \type{empty} parameter:

\startbuffer
\setupquestion[showanswer=false,point=2,showpoint=false]

\startquestion[showanswer=false]
  Different fill-in style examples:
  \startproblem
    {\setupfillin[empty=number,n=5]
     Numbered style: Lorem ipsum \fillin{answer}}
    
    {\setupfillin[empty=yes,n=5]
     Hidden style: Lorem ipsum \fillin{answer}}
    
    {\setupfillin[empty=none,n=5]
     Display style: Lorem ipsum \fillin{answer}}
  \stopproblem
\stopquestion
\stopbuffer

\typebuffer \getbuffer

\subsection{True/False Questions}

Fill-in-the-blank can also be used for true/false questions:

\startbuffer
\setupfillin[empty=none,style=\tta]
\startquestion
Determine whether the following statements are correct.
  \startproblem
  \startpitem \fillin{T} This is a correct statement.\stoppitem
  \startpitem \fillin{F} This is an incorrect statement.\stoppitem
  \startpitem This is a correct statement.\hfill\fillin{\symbol{markcheck}}\stoppitem
  \startpitem This is an incorrect statement.\hfill\fillin{\symbol{markcross}}\stoppitem
  \stopproblem
\stopquestion
\stopbuffer

\typebuffer \getbuffer

\chapter{Reading Material}

\section{Basic Usage}

Reading comprehension is a common question type in exams, usually containing a reading passage and several related questions.
\type{material} environment is used to place reading material, supporting title, author, source, and other information.

\startbuffer
\setupmaterial[title][style=\ssa]
\setupmaterial[author][style=\tta]
\setupmaterial[source][style=ss]

\startmaterial[title={Material Title},author={Author Name},source={Source}]
  This is the content of the reading material.
  It can contain multiple paragraphs.
  \indicator{content to be questioned one}
  \indicator{content to be questioned two}
\stopmaterial

\startquestion
  \indicator{content to be questioned one}
  \choice{Option A,Option B,Option C,Option D}
\stopquestion

\startquestion
  \indicator{content to be questioned two}
  \choice{Option A,Option B,Option C,Option D}
\stopquestion
\stopbuffer

\typebuffer \getbuffer

\section{Parameter Description}

\showsetup[setupmaterial]

\starttabulate[|l|l|p|]
\HL
\NC Parameter \NC Default \NC Description \NC\NR
\HL
\NC title \NC \NC Material title \NC\NR
\NC author \NC \NC Author \NC\NR
\NC source \NC \NC Source \NC\NR
\NC style \NC \NC Body font style \NC\NR
\NC color \NC \NC Body font color \NC\NR
\NC before \NC \NC Content before material \NC\NR
\NC after \NC \NC Content after material \NC\NR
\HL
\stoptabulate

\section{Advanced Usage}

\subsection{indicator Command}

\type{\indicator} command is used to mark question points in the material, establishing the relationship between material and questions.
Marked content is displayed with a special style in the material and referenced in the questions.

\startitemize[packed]
\item Use \type{\indicator{marker name}} in material to mark key content
\item Use \type{\indicator{marker name}} in questions to reference that content
\item Association is automatically established during compilation
\stopitemize

\startbuffer
\startmaterial[title={Example Material}]
  Article content, where \indicator{key content one} needs to be questioned,
  and \indicator{key content two} also needs to be questioned.
\stopmaterial
\stopbuffer

\typebuffer \getbuffer

\chapter{Verse Typesetting}

\section{Basic Usage}

\type{verse} environment is used to typeset poetry, lyrics, and other verse content:

\startbuffer
\startverse[title=Linjiang Xian $\cdot$ Rolling Yangtze River Eastward,author=Ming $\cdot$ Yang Shen,sample={White-haired fisherman on the river isle}]
滾滾長江東逝水浪花淘盡英雄\\
是非成敗轉頭空\\
青山依舊在 幾度夕陽紅\\
白髮漁樵江渚上 慣看秋月春風\\
一壺濁酒喜相逢\\
古今多少事 都付笑談中
\stopverse
\stopbuffer

\typebuffer \getbuffer

\section{Parameter Description}

\showsetup[setupverse]

\starttabulate[|l|l|p|]
\HL
\NC Parameter \NC Default \NC Description \NC\NR
\HL
\NC title \NC \NC Poem title \NC\NR
\NC author \NC \NC Author information \NC\NR
\NC source \NC \NC Source information \NC\NR
\NC sample \NC \NC Sample text for calculating indentation \NC\NR
\NC skip \NC both \NC Indentation mode: left, right, both, none \NC\NR
\NC align \NC \NC Alignment \NC\NR
\NC style \NC \NC Font style \NC\NR
\NC color \NC \NC Text color \NC\NR
\NC before \NC \NC Content before environment \NC\NR
\NC after \NC \NC Content after environment \NC\NR
\HL
\stoptabulate

\section{Advanced Usage}

Styles for title, author, and source can be set separately:

\startbuffer
\setupverse[sample={滾滾長江東逝水浪花淘盡英雄},skip=left,before=\blank,after=\blank]
\setupverse[title][align=center,style=bf,skip=both]
\setupverse[author][align=flushright,style=it,skip=both]
\setupverse[source][align=flushright,style=it,skip=both]
\stopbuffer

\typebuffer

\chapter{Cloze Test}

\section{Basic Usage}

\type{close} environment is used to create cloze tests, embedding options in text for students to select.
Each option is marked using \type{\closechoice} command, correct answers are marked with \type{[*]}.

\startbuffer
\startclose[showanswer=true,showpoint=true,point=1.5,answer=nil]
On Oct. 11, hundreds of runners competed in a cross-country race.
Melanie Bailey should have \closechoice[designed,followed,changed,finished] the
course earlier than she did. Her \closechoice[delay,chance,trouble,excuse] came
because she was carrying a \closechoice[judge,volunteer,classmate,competitor]
across the finish line.
\stopclose
\stopbuffer

\typebuffer \getbuffer

\section{Parameter Description}

\starttabulate[|l|l|p|]
\HL
\NC Parameter \NC Default \NC Description \NC\NR
\HL
\NC reset \NC false \NC Whether to reset question number \NC\NR
\NC n \NC 8 \NC Underline length \NC\NR
\NC width \NC 1.5em \NC Question number width \NC\NR
\NC answerwidth \NC 4em \NC Answer area width \NC\NR
\NC point \NC \NC Points per question \NC\NR
\NC answer \NC \NC Default answer \NC\NR
\NC showanswer \NC \NC Whether to show answer \NC\NR
\NC showpoint \NC \NC Whether to show score \NC\NR
\NC style \NC \NC Text style \NC\NR
\NC color \NC \NC Text color \NC\NR
\NC barstyle \NC \NC Underline style \NC\NR
\NC barcolor \NC \NC Underline color \NC\NR
\NC frame \NC \NC Whether to show frame \NC\NR
\NC stopper \NC \NC Question number suffix \NC\NR
\NC distance \NC \NC Option spacing \NC\NR
\NC before \NC \NC Content before environment \NC\NR
\NC after \NC \NC Content after environment \NC\NR
\HL
\stoptabulate

\section{Advanced Usage}

\subsection{closechoice Command}

\type{\closechoice} command is used to insert options in cloze tests, syntax:
\type{\closechoice[option1,option2,option3,...]}.

Correct answers are marked with \type{[*]}: \type{\closechoice[option1,{[*]correct option},option3]}.

\subsection{optclose Environment}

\type{optclose} environment is a variant of \type{close}, used for option-based fill-in:

\startbuffer
\startoptclose[point=2,showanswer=true]
The latest \ocitem[answer=engineering]{engineer} techniques are applied to create
this protective \ocitem[answer=functional]{function} structure.
\stopoptclose
\stopbuffer

\typebuffer \getbuffer

\chapter{Writing Area}

\section{Basic Usage}

\subsection{Composition Grid}

Use \type{\writingbox} to create composition grids:

\startbuffer
\writingbox
  [column=6,
   row=6,
   evencolor=blue,
   oddcolor=darkgray,
   evenheight=1cm,
   oddheight=.5cm,
   height=10cm,
   width=\textwidth]{}
\stopbuffer

\typebuffer \getbuffer

\subsection{Grid Paper}

Use \type{\writinggrid} to create grid paper:

\startbuffer
\writinggrid
  [cellwidth=1cm,
   column=12,
   row=8,
   cellcolor=darkgray,
   height=8cm,
   width=\textwidth]{}
\stopbuffer

\typebuffer \getbuffer

\section{Parameter Description}

\subsection{writingbox Parameters}

\showsetup[setupwritingbox]

\starttabulate[|l|l|p|]
\HL
\NC Parameter \NC Default \NC Description \NC\NR
\HL
\NC column \NC \NC Number of columns \NC\NR
\NC row \NC \NC Number of rows \NC\NR
\NC evencolor \NC \NC Even row color \NC\NR
\NC oddcolor \NC \NC Odd row color \NC\NR
\NC evenheight \NC \NC Even row height \NC\NR
\NC oddheight \NC \NC Odd row height \NC\NR
\NC height \NC \NC Total height \NC\NR
\NC width \NC \NC Total width \NC\NR
\NC align \NC \NC Alignment \NC\NR
\NC rulethickness \NC \NC Line thickness \NC\NR
\HL
\stoptabulate

\subsection{writinggrid Parameters}

\showsetup[setupwritinggrid]

\starttabulate[|l|l|p|]
\HL
\NC Parameter \NC Default \NC Description \NC\NR
\HL
\NC column \NC \NC Number of columns \NC\NR
\NC row \NC \NC Number of rows \NC\NR
\NC cellwidth \NC \NC Cell width \NC\NR
\NC cellcolor \NC \NC Cell color \NC\NR
\NC height \NC \NC Total height \NC\NR
\NC width \NC \NC Total width \NC\NR
\HL
\stoptabulate

\section{Advanced Usage}

\subsection{writing Environment}

\type{writing} environment is used to create composition questions, supporting score and answer settings:

\startbuffer
\startwriting[point=45,showpoint=true]
  Composition prompt
  
  Please write according to the following requirements:
  \startitemize
  \item Requirement one
  \item Requirement two
  \stopitemize
  
  \startanswer
  Sample essay content...
  \stopanswer
\stopwriting
\stopbuffer

\typebuffer \getbuffer

\subsection{subwriting Environment}

\type{subwriting} environment is used to create sub-composition questions:

\startbuffer
\startwriting[point=100]
  Composition prompt
  
  \startsubwriting[point=50]
    Sub-composition one
    \startanswer
    Sample essay one...
    \stopanswer
  \stopsubwriting
  
  \startsubwriting[point=50]
    Sub-composition two
    \startanswer
    Sample essay two...
    \stopanswer
  \stopsubwriting
\stopwriting
\stopbuffer

\typebuffer \getbuffer

\subsection{writing Parameters}

\type{writing} environment inherits from \type{question} environment, parameters are set through \type{\setupwriting}:

\starttabulate[|l|l|p|]
\HL
\NC Parameter \NC Default \NC Description \NC\NR
\HL
\NC point \NC \NC Composition score \NC\NR
\NC answer \NC See sample essay \NC Reference answer \NC\NR
\NC totalanswer \NC \NC Total reference answer \NC\NR
\NC showanswer \NC \NC Whether to show answer \NC\NR
\NC showpoint \NC \NC Whether to show score \NC\NR
\NC style \NC \NC Question font style \NC\NR
\NC color \NC \NC Question color \NC\NR
\NC option \NC n,packed \NC List options \NC\NR
\NC before \NC \NC Content before environment \NC\NR
\NC after \NC \NC Content after environment \NC\NR
\HL
\stoptabulate

\chapter{Symbol Definitions}

The module has built-in various marker symbols (requires NotoSansSymbols2 font):

\starttabulate[|l|r|l|r|]
\HL
\NC Command \NC Symbol \NC Command \NC Symbol \NC\NR
\HL
\NC \type{\symbol{marksquare}} \NC \symbol{marksquare}
\NC \type{\symbol{markcircle}} \NC \symbol{markcircle} \NC\NR
\NC \type{\symbol{markcheck}} \NC \symbol{markcheck}
\NC \type{\symbol{markcross}} \NC \symbol{markcross} \NC\NR
\NC \type{\symbol{markheavycheck}} \NC \symbol{markheavycheck}
\NC \type{\symbol{markheavycross}} \NC \symbol{markheavycross} \NC\NR
\NC \type{\symbol{marksquarecheck}} \NC \symbol{marksquarecheck}
\NC \type{\symbol{marksquarecross}} \NC \symbol{marksquarecross} \NC\NR
\NC \type{\symbol{markcirclecheck}} \NC \symbol{markcirclecheck}
\NC \type{\symbol{markcirclecross}} \NC \symbol{markcirclecross} \NC\NR
\HL
\stoptabulate

\section{Circled Text}

\type{\circledtext} command is used to create text with circles or boxes:

\startbuffer
\circledtext{1}\quad
\circledtext[radius=0.5]{A}\quad
\circledtext[radius=0,fillcolor=black,color=white]{B}\quad
\circledtext[radius=0.2]{C}
\stopbuffer

\typebuffer \getbuffer

\subsection{circledtext Parameters}

\showsetup[setupcircledtext]

\starttabulate[|l|l|p|]
\HL
\NC Parameter \NC Default \NC Description \NC\NR
\HL
\NC frametype \NC circle \NC Frame type: circle, square \NC\NR
\NC framecolor \NC \NC Frame color \NC\NR
\NC fillcolor \NC \NC Fill color \NC\NR
\NC dashcolor \NC \NC Dashed line color \NC\NR
\NC dashlinewidth \NC \NC Dashed line width \NC\NR
\NC color \NC \NC Text color \NC\NR
\NC radius \NC \NC Corner radius \NC\NR
\NC effect \NC normal \NC Effect \NC\NR
\NC scaled \NC \NC Scale ratio \NC\NR
\HL
\stoptabulate

\subsection{Circled Number Conversion}

The module provides various circled number conversion formats for list numbering:

\startbuffer
\startitemize[Co:num,packed]
  \item First item
  \item Second item
  \item Third item
\stopitemize

\startitemize[CO:num,packed]
  \item First item
  \item Second item
  \item Third item
\stopitemize
\stopbuffer

\typebuffer \getbuffer

\starttabulate[|l|l|]
\HL
\NC Conversion Format \NC Description \NC\NR
\HL
\NC Co:num \NC White background black text circled number \NC\NR
\NC CO:num \NC Black background white text circled number \NC\NR
\NC Cr:num \NC White background black text rounded box number \NC\NR
\NC CR:num \NC Black background white text rounded box number \NC\NR
\NC Cs:num \NC White background black text square box number \NC\NR
\NC CS:num \NC Black background white text square box number \NC\NR
\NC Co:Alph \NC Uppercase circled letter \NC\NR
\NC Co:alph \NC Lowercase circled letter \NC\NR
\NC Co:cn \NC Circled Chinese number \NC\NR
\NC Co:hira \NC Circled hiragana \NC\NR
\NC Co:kata \NC Circled katakana \NC\NR
\HL
\stoptabulate

\chapter{Dictionary Entries}

\section{Basic Usage}

\type{lexicalentry} environment is used to create dictionary-style entries, suitable for making vocabulary lists, glossaries, or language learning materials.
Each entry can contain part of speech tags, accent marks, definitions, example sentences, etc.

\startbuffer
\startlexicalentry[class=n.,accent=1]{太陽}
  \startlexicallists
    \item 太陽系の中心をなし、地球にもっとも近い恒星。お日さま。\antonym{太陰}。
    \item 光明・希望の象徴。「心の\alias」
  \stoplexicallists
\stoplexicalentry
\stopbuffer

\typebuffer \getbuffer

\section{Parameter Description}

\showsetup[setuplexicalentry]

\starttabulate[|l|l|p|]
\HL
\NC Parameter \NC Default \NC Description \NC\NR
\HL
\NC class \NC \NC Part of speech tag (e.g., "n.", "v.") \NC\NR
\NC accent \NC \NC Accent mark (number) \NC\NR
\NC number \NC \NC Number \NC\NR
\NC style \NC \NC Entry style \NC\NR
\NC color \NC \NC Entry color \NC\NR
\NC before \NC \NC Content before entry \NC\NR
\NC after \NC \NC Content after entry \NC\NR
\HL
\stoptabulate

\section{Advanced Usage}

\subsection{Dictionary Lists}

\type{lexicallists} is a supporting list environment, supporting four levels of nesting, each using different numbering styles:

\startbuffer
\startlexicalentry[class=v.]{走る}
  \startlexicallists
    \item 足を動かして進む。\antonym{止まる}。
    \item 逃げる。
      \startlexicallists
        \item 時間が経過する。
        \item 機械が動く。
      \stoplexicallists
  \stoplexicallists
\stoplexicalentry
\stopbuffer

\typebuffer \getbuffer

\subsection{Dictionary Enumerations}

The module provides various dictionary enumeration environments for structured display of entry information:

\starttabulate[|l|l|]
\HL
\NC Environment \NC Description \NC\NR
\HL
\NC lexicalgrammar \NC Grammar explanation \NC\NR
\NC lexicalmeaning \NC Definition explanation \NC\NR
\NC lexicaljunction \NC Conjunction explanation \NC\NR
\NC lexicalexample \NC Example explanation \NC\NR
\NC lexicalcomparing \NC Comparison explanation \NC\NR
\NC lexicalsummary \NC Summary explanation \NC\NR
\HL
\stoptabulate

\startbuffer
\startlexicalentry[class=adj.]{美しい}
  \startlexicalmeaning
    見た目や心がすばらしい。きれいだ。
  \stoplexicalmeaning
  
  \startlexicalexample
    美しい景色を眺める。\\
    美しい心を持つ。
  \stoplexicalexample
  
  \startlexicalcomparing
    「美しい」は芸術的な美しさ、「きれい」は清潔さを強調。
  \stoplexicalcomparing
\stoplexicalentry
\stopbuffer

\typebuffer \getbuffer

\subsection{Association Commands}

Dictionary entries support various association commands for marking relationships between entries:

\starttabulate[|l|p(8cm)|]
\HL
\NC Command \NC Description \NC\NR
\HL
\NC \type{\antonym{word}} \NC Antonym \NC\NR
\NC \type{\synonym{word}} \NC Synonym \NC\NR
\NC \type{\alias} \NC Entry alias reference \NC\NR
\NC \type{\refer{word}} \NC See also reference \NC\NR
\NC \type{\example{example}} \NC Example mark \NC\NR
\NC \type{\junction{junction}} \NC Junction mark \NC\NR
\HL
\stoptabulate

\chapter{Advanced Features}

\section{Unlimited Nesting Support}

The module supports unlimited levels of question nesting with automatic answer tree management:

\startbuffer
\startquestion[point=20]
First level major question
  \startquestion[point=10]
  Second level sub-question
    \startquestion[point=5]
    Third level sub-question
    \stopquestion
  \stopquestion
\stopquestion
\stopbuffer

\typebuffer \getbuffer

\section{Data Persistence}

Answer data is automatically saved to .tuc file and can be accessed during second compilation:

\starttyping
% Access saved answer data
\datasetvariable{UnlimitedDataset}{question ID}{answer}
\datasetvariable{UnlimitedDataset}{question ID}{point}
\stoptyping

\section{Answer Tree Structure}

The basicexam module uses an answer tree structure to manage all question answers and score information.
Each question node contains the following attributes:

\starttabulate[|l|p(6cm)|l|]
\HL
\NC Attribute \NC Description \NC Retrieval Method \NC\NR
\HL
\NC id \NC Node unique identifier (e.g., \type{jobname-1-Q-1}) \NC \type{\UnlimitedCurrentID} \NC\NR
\NC order \NC Level order (e.g., \type{1}, \type{2-1}, \type{2-2-1}) \NC \type{\UnlimitedGetOrder{ID}} \NC\NR
\NC answer \NC Answer content \NC \type{\UnlimitedDisplayAnswer{ID}} \NC\NR
\NC point \NC Current node score (valid for leaf nodes) \NC \type{\UnlimitedDisplayPoint{ID}} \NC\NR
\NC totalpoint \NC Total score of this node and child nodes \NC \type{\UnlimitedTotalPoint{ID}} \NC\NR
\NC totalnumber \NC Total count of this node and child nodes \NC \type{\UnlimitedTotalNumber{ID}} \NC\NR
\NC is_leaf \NC Whether it is a leaf node \NC \type{\UnlimitedIsLeaf{ID}} \NC\NR
\NC depth \NC Node depth (0=root node) \NC \type{\UnlimitedGetDepth} \NC\NR
\HL
\stoptabulate

\section{Node Information Retrieval}

\subsection{Current Node Information}

\starttabulate[|l|p(8cm)|]
\HL
\NC Command \NC Description \NC\NR
\HL
\NC \type{\UnlimitedCurrentID} \NC Get current node ID \NC\NR
\NC \type{\UnlimitedGetDepth} \NC Get current node depth \NC\NR
\NC \type{\UnlimitedCurrentAnswer} \NC Get current node answer \NC\NR
\NC \type{\UnlimitedCurrentPoint} \NC Get current node total score \NC\NR
\HL
\stoptabulate

\subsection{Specified Node Information}

Use the following commands to get attributes of specified nodes:

\startitemize[packed]
\item \type{\UnlimitedDisplayPoint{ID}} - Get node score
\item \type{\UnlimitedDisplayAnswer{ID}} - Get node answer
\item \type{\UnlimitedTotalPoint{ID}} - Get node total score (including sub-questions)
\item \type{\UnlimitedTotalNumber{ID}} - Get node total count (including sub-questions)
\item \type{\UnlimitedGetOrder{ID}} - Get node level order
\item \type{\UnlimitedIsLeaf{ID}} - Check if it is a leaf node
\stopitemize

\startbuffer
% Get score of question 1
\UnlimitedDisplayPoint{jobname-1-Q-1}

% Get answer of question 1
\UnlimitedDisplayAnswer{jobname-1-Q-1}

% Get total score of question 1 (including sub-questions)
\UnlimitedTotalPoint{jobname-1-Q-1}

% Get total count of question 1 (including sub-questions)
\UnlimitedTotalNumber{jobname-1-Q-1}
\stopbuffer

\typebuffer

\section{Structure Level Statistics}

The module provides statistics by chapter/section level, allowing you to get question counts and scores for specified levels:

\starttabulate[|l|p(8cm)|]
\HL
\NC Command \NC Description \NC\NR
\HL
\NC \type{\UnlimitedStructureTotalNumber{level}{number}} \NC Top-level question count for specified level (major questions) \NC\NR
\NC \type{\UnlimitedStructureTotalPoint{level}{number}} \NC Total points for specified level \NC\NR
\NC \type{\UnlimitedStructureLeafNumber{level}{number}} \NC Top-level leaf node count for specified level (independent major questions) \NC\NR
\NC \type{\UnlimitedStructureAllLeafNumber{level}{number}} \NC All leaf node count for specified level (actual sub-questions) \NC\NR
\NC \type{\UnlimitedCurrentStructureTotalNumber{level}} \NC Top-level question count for current level \NC\NR
\NC \type{\UnlimitedCurrentStructureTotalPoint{level}} \NC Total points for current level \NC\NR
\NC \type{\UnlimitedCurrentStructureLeafNumber{level}} \NC Top-level leaf node count for current level \NC\NR
\NC \type{\UnlimitedCurrentStructureAllLeafNumber{level}} \NC All leaf node count for current level \NC\NR
\HL
\stoptabulate

\subsection{Parameter Description}

\startitemize[packed]
\item \type{level}: Level number, 1=chapter, 2=section, 3=subsection...
\item \type{number}: The number of that level (e.g., chapter 1, chapter 2)
\stopitemize

{\bf Important}: When using \type{\UnlimitedCurrentStructure*} commands, they must be called within the corresponding chapter/section structure, otherwise they will return 0.

\subsection{Statistics Concepts}

\startitemize[packed]
\item {\bf Major questions}: Top-level questions directly under chapters/sections, not contained within other questions
\item {\bf Independent major questions}: Top-level questions without sub-questions, i.e., standalone major questions
\item {\bf All sub-questions}: All actual questions that need to be answered, including nested sub-questions
\stopitemize

\subsection{Statistics Example}

Assume Chapter 1 has the following question structure:

\startbuffer
Chapter 1
├── Question 1 (has sub-questions)
│   ├── Sub-question 1.1
│   └── Sub-question 1.2
├── Question 2 (standalone question)
└── Question 3 (has sub-questions)
    ├── Sub-question 3.1
    ├── Sub-question 3.2
    └── Sub-question 3.3
\stopbuffer

\typebuffer

Statistics results:
\startitemize[packed]
\item Major questions: 3 (Question 1, Question 2, Question 3)
\item Independent major questions: 1 (only Question 2 is standalone)
\item All sub-questions: 6 (Sub-question 1.1, 1.2, Question 2, Sub-question 3.1, 3.2, 3.3)
\stopitemize

\section{Question Level Statistics}

In addition to chapter-level statistics, the module also provides question-level hierarchy statistics:

\starttabulate[|l|p(8cm)|]
\HL
\NC Command \NC Description \NC\NR
\HL
\NC \type{\UnlimitedQuestionLevelNumber{ID}{depth}} \NC Get the number of sub-questions at a specific level for a question \NC\NR
\NC \type{\UnlimitedQuestionLeafNumber{ID}} \NC Get the total leaf node count for a question (sub-questions) \NC\NR
\HL
\stoptabulate

\subsection{Parameter Description}

\startitemize[packed]
\item \type{ID}: Question node ID (e.g., \type{jobname-1-Q-1})
\item \type{depth}: Relative depth level (1=first level sub-questions, 2=second level sub-questions...)
\stopitemize

\section{Extra Data}

Additional node information can be set through \type{\UnlimitedSetExtraData{key}{value}}, which can be retrieved via \type{\UnlimitedGetCurrentExtraData{key}}.

Common extra data includes:

\starttabulate[|l|p(8cm)|]
\HL
\NC Key \NC Description \NC\NR
\HL
\NC totalitemnumber \NC Number of options in choice environment \NC\NR
\NC availablecolumns \NC Available columns \NC\NR
\NC itemmaxwidth \NC Maximum width of options \NC\NR
\HL
\stoptabulate

\section{Usage Examples}

\subsection{Getting Chapter Statistics}

\startbuffer
% Get statistics for chapter 1
Chapter 1 question count: \UnlimitedStructureTotalNumber{1}{1}
Chapter 1 total score: \UnlimitedStructureTotalPoint{1}{1}
Chapter 1 leaf node count: \UnlimitedStructureLeafNumber{1}{1}
\stopbuffer

\typebuffer

\subsection{Getting Question Information}

\startbuffer
% Get information for a specific question
Question 1 score: \UnlimitedGetAnswer{jobname-1-Q-1}{point}
Question 1 answer: \UnlimitedGetAnswer{jobname-1-Q-1}{answer}
Question 1 total score: \UnlimitedGetAnswer{jobname-1-Q-1}{totalpoint}
Question 1 total count: \UnlimitedGetAnswer{jobname-1-Q-1}{totalnumber}
\stopbuffer

\typebuffer

\subsection{Getting Information Inside Questions}

\startbuffer
\question[point=5,answer=A]{Question
  Current ID: \UnlimitedCurrentID
  Current depth: \UnlimitedGetDepth
  Current answer: \UnlimitedDisplayAnswer{\UnlimitedCurrentID}
  Current score: \UnlimitedDisplayPoint{\UnlimitedCurrentID}
}
\stopbuffer

\typebuffer

\section{Debug Features}

In check mode, the module automatically outputs debug information:

\starttyping
\usemodule[basicexam][mode={teacher,check}]
\stoptyping

You can also manually call debug commands:

\starttyping
% Output answer tree structure
\UnlimitedDebug

% Export answers to TSV file
\UnlimitedExport{answers.tsv}
\stoptyping

\chapter{Contributing and Modifying}

\section{Submitting Requests}

If you find bugs or have new feature requests during use, you can submit them through:

\startitemize[packed]
\item Submit an Issue on GitHub repository
\item Send email to author
\item Post on relevant forums for discussion
\stopitemize

When submitting requests, please try to provide the following information:

\startitemize[packed]
\item ConTeXt version information (run \type{context --version})
\item Minimal reproduction example
\item Expected behavior and actual behavior
\stopitemize

\section{Modifying This Module}

This module is written in LuaMetaTeX, source code is located in \type{t-basicexam.mklx} file.

\subsection{Module Structure}

\starttabulate[|l|p(8cm)|]
\HL
\NC File \NC Description \NC\NR
\HL
\NC t-basicexam.mklx \NC Main module file, contains all command definitions \NC\NR
\NC i-basicexam.xml \NC Interface definition file, for command documentation generation \NC\NR
\NC doc/basicexam-manual.tex \NC User manual \NC\NR
\NC doc/test-exam.tex \NC Test file \NC\NR
\HL
\stoptabulate

\subsection{Modification Suggestions}

\startitemize[packed]
\item Please backup original files before modifying
\item Recommended to read source code comments first
\item Please run test files for verification after modification
\item For major improvements, Pull Requests are welcome
\stopitemize

\stoptext
